% Append-only technical-record section. Requires amsmath and booktabs.
\section{Exploratory Edge-Layer Mixed-Precision Experiment}
\label{sec:qwen30b-q2q3-edge}

\subsection{Predeclared question and scope}

Following promotion of the resident Qwen3-30B-A3B selective-Q2 expert
candidate, an explicitly quality-gated experiment asked whether allocating
higher expert precision to the first and last eight transformer layers could
recover quality while retaining at least 50 generated tokens per second on the
18~GB M3 Pro.  This was a changed-weight experiment, not an exact-output
optimisation.  The edge-layer allocation was declared as an architectural
heuristic and was not represented as a measured sensitivity result.

The experiment plan was written before candidate creation and is stored as
\texttt{qwen3-30b-a3b-q2q3-edge-plan-v1.json}, with SHA-256
\texttt{4b6912d12ef983e71dbe0db307095758541c7452522560fe3640e0a47eb12542}.
The frozen acceptance rules required a structured strict score above 6/8, no
semantic regression below 7/8, no multiple-choice regression below 11/12,
repeatable parseable output, and replicated 256-token performance at or above
the declared speed thresholds.

\subsection{Candidate construction and physical boundary}

Expert tensors in layers 0--7 and 40--47 were quantised to Q3\_K.  Expert
tensors in layers 8--39 remained Q2\_K, while non-expert tensors retained the
Q4\_K\_M base policy.  The resulting candidate is 11,497,672,192 bytes
(10,959.34~MiB) with SHA-256
\texttt{1cec6825acb0a912a449ac0f5aacfc376ca4dce548d1317b741a3253438287f4}.
It was built on temporary internal storage because the SD evidence volume had
only approximately 9~GiB free.  Consequently this experiment makes no claim
about SD cold-load speed.  Neither SD-resident source model was overwritten or
deleted.

The candidate loaded successfully with all model layers on Metal.  A narrow
32-prompt-token, 16-generation-token smoke test measured 164.14 prompt
tokens/s and 58.13 generation tokens/s.  This established executability, not
the predeclared replicated long-generation performance result.

\subsection{Frozen quality result}

\begin{table}[ht]
\centering
\begin{tabular}{lccc}
\toprule
Evaluation & Q2 control & Q2/Q3 edge candidate & Gate \\
\midrule
Structured strict & 6/8 & 7/8 & pass \\
Structured semantic & 7/8 & 8/8 & pass \\
Structured parseable & 8/8 & 8/8 & pass \\
Multiple choice & 11/12 & 10/12 & fail \\
\bottomrule
\end{tabular}
\caption{Frozen quality results for the heuristic mixed-precision candidate.}
\end{table}

The structured cohort was run twice.  Both raw text and parsed values were
identical between runs.  The sole strict miss returned
\texttt{unit=``gram''} rather than the schema-exact \texttt{unit=``g''}, while
the numerical conversion was correct.  In the multiple-choice cohort, the
candidate missed cases \texttt{m03} and \texttt{m04}, producing 10/12 versus
the Q2 control's 11/12.  Because the predeclared multiple-choice gate failed,
the replicated 256-token performance stage was not run.

\subsection{Decision and learning}

The candidate is \textbf{NOT\_PROMOTED}.  Restoring precision at model edges
improved one frozen structured cohort but regressed the separate
multiple-choice cohort.  The result falsifies a general claim that first/last
layers are automatically the best precision allocation for this model.  Any
subsequent mixed-precision candidate must be based on task-diverse, measured
layer-sensitivity evidence rather than positional intuition alone.

The canonical result is
\texttt{qwen3-30b-a3b-q2q3-edge-result-v1.json}, SHA-256
\texttt{d47d6e0c65522c31f6c290c31b4cdba342c52d4c7a97a7860a7cf9d483eb0475}.
